\section{A weighted matching consequence of hypergraph edge colouring}
\label{sec:rhost-weighted-matching}

We use the Pippenger--Spencer edge-colouring theorem:
for fixed uniformity, an almost regular hypergraph with relative codegree $o(1)$ has chromatic index $(1+o(1))$ times its degree.
The primary paper is
N.~Pippenger and J.~Spencer, \emph{Asymptotic behavior of the chromatic index for hypergraphs},
Journal of Combinatorial Theory, Series~A \textbf{51} (1989), 24--42,
\url{https://doi.org/10.1016/0097-3165(89)90074-5}.
A precise statement also appears as Lemma~2.1 of Alon and Yuster,
\url{https://web.math.princeton.edu/~nalon/PDFS/ayhyper2.pdf}.

\begin{lemma}[Maximum degree and a weighted objective]\label{lem:rhost-max-degree-colouring}
Fix $r\ge3$.
Let $H$ be a simple $r$-uniform hypergraph with maximum degree at most $D\to\infty$ and maximum codegree $o(D)$.
Give each edge a nonnegative value $v_e$.
Then $H$ has a matching $M$ such that
\[
 \sum_{e\in M}v_e\ge\frac{1-o(1)}D\sum_{e\in E(H)}v_e.
\]
No lower bound on vertex degrees is required.
\end{lemma}

\begin{proof}
We reduce to the regular form of the edge-colouring theorem.
Take $L$ disjoint copies of $H$, where $L$ is a sufficiently large multiple of $r$.
For each original vertex $u$, group its $L$ copies together.
Independently choose $D-\deg_H(u)$ uniformly random partitions of that group into $r$-sets.
Add all sets in those partitions as edges.
Every vertex now has degree $D$, counted with multiplicity.

For a fixed pair within one group, its probability of sharing a part in one partition is $(r-1)/(L-1)$.
The probability that it shares a part in at least three partitions is $O(D^3/L^3)$.
A union bound over all groups and pairs gives $O(|V(H)|D^3/L)$.
The expected number of repeated added edges is $O(|V(H)|D^2L^{2-r})$.
Choose $L$ so large that both bounds tend to zero.
There is therefore a choice with no repeated added edges and with every new pair codegree at most two.

An original edge lies within one copy and uses distinct vertex positions.
An added edge lies within one vertex-position group and uses distinct copies.
No pair is contained in both an original and an added edge.
The resulting simple hypergraph is $D$-regular and has codegree at most $\max\{\Delta_2(H),2\}=o(D)$.
By Pippenger--Spencer, its edges split into $(1+o(1))D$ matchings.
Give added edges value zero and copied original edges their original values.
Some colour class has total value at least
\[
 \frac{L}{(1+o(1))D}\sum_{e\in E(H)}v_e.
\]
Within that colour class, some copy of $H$ carries at least the average of this value over the $L$ copies.
Its original edges form the required matching.
\end{proof}

\begin{lemma}[Weighted length packing]\label{lem:rhost-weighted-matching}
Fix integers $r\ge3$ and $C\ge1$, and constants $a,V,c>0$.
For $h\to\infty$, let $\mathcal A_h$ be a family of distinct $r$-element sets on at most $h^C$ vertices.
Give each set a nonnegative weight $p_A$ and a value $v_A\in[0,V]$.
Suppose
\[
 \sum_Ap_A=O(1),\qquad \sum_Ap_Av_A\ge c,
\]
\[
 \max_z\sum_{A\ni z}p_A\le\frac{a+o(1)}h,
 \qquad
 \max_{z\ne z'}\sum_{A\supset\{z,z'\}}p_A=O(h^{-3/2}).
\]
Then some pairwise disjoint subfamily $\mathcal M_h$ satisfies
\[
 \sum_{A\in\mathcal M_h}v_A
 \ge\left(\frac h a-o(h)\right)\sum_Ap_Av_A.
\]
\end{lemma}

\begin{proof}
Put $\lambda=h(\ln h)^2$.
Each $p_A$ is bounded by a pair weight, so $\lambda p_A<1$ for large $h$.
Include every distinct set independently with probability $\lambda p_A$.
This gives a simple $r$-uniform hypergraph $H$.

Every vertex degree has mean at most $(a+o(1))(\ln h)^2$.
Chernoff bounds, followed by a union bound over polynomially many vertices, give
\[
 \Delta(H)\le(a+o(1))(\ln h)^2
\]
with probability tending to one.
For example, an additional relative allowance $(\ln h)^{-1/4}$ suffices.
Every codegree has mean $O(h^{-1/2}(\ln h)^2)$.
A Chernoff upper-tail bound at $K\ln h$, followed by a union bound over polynomially many pairs, gives
$\Delta_2(H)=O(\ln h)$.

Let $W=\sum_{A\in E(H)}v_A$ and $S=\sum_Ap_Av_A$.
Then $\mathbb EW=\lambda S$ and
\[
 \operatorname{Var}(W)\le V^2\lambda\sum_Ap_A=O(\lambda).
\]
Since $S\ge c$, Chebyshev's inequality gives $W=(1-o(1))\lambda S$ with probability tending to one.
Fix a realization satisfying all these conclusions.
If needed, increase the degree parameter to an integer
$D=(a+o(1))(\ln h)^2\ge\Delta(H)$ in the regularisation proof of Lemma~\ref{lem:rhost-max-degree-colouring}.
That proof remains valid with this upper degree bound.
Its matching has value at least
\[
 \frac{(1-o(1))\lambda S}{(a+o(1))(\ln h)^2}
 =\left(\frac h a-o(h)\right)S.
\]
\end{proof}

\paragraph{Variable set sizes and private vertices.}
The lemma also applies when every real set has size between two and a fixed integer $r$.
Aggregate equal real sets before padding, retaining their weight and a fixed objective value for that set.
If a real set $A$ has fewer than $r$ vertices, add $r-|A|$ new private vertices used only by that set.
Every individual weight $p_A$ is at most a weighted codegree of two real vertices of $A$.
Thus each private vertex has weighted degree $O(h^{-3/2})$.
Every new pair codegree has the same bound, including pairs of private vertices from one padded set.
If the real ground set has $O(h)$ vertices, there are polynomially many possible real sets and private vertices.
The hypotheses therefore remain valid for the padded $r$-uniform family.
A matching of padded sets restricts to disjoint real sets and preserves its objective value.

In particular, wheel outcomes may retain $k\in\{4,\ldots,12\}$ required non-rim lengths and one marked rim length.
Their real sets have sizes $k+1\in\{5,\ldots,13\}$.
After padding to uniformity thirteen, their values $k-3\in\{1,\ldots,9\}$ count the saved block rank.
